\documentclass[11pt,a4paper]{article}
\usepackage[margin=2.6cm]{geometry}
\usepackage[T1]{fontenc}
\usepackage{lmodern}
\usepackage{microtype}
\usepackage{booktabs}
\usepackage{array}
\usepackage{enumitem}
\usepackage{xcolor}
\usepackage{fancyvrb}
\usepackage[hidelinks]{hyperref}

\definecolor{codebg}{gray}{0.95}
\DefineVerbatimEnvironment{code}{Verbatim}{fontsize=\small,frame=single,framerule=0.2pt,framesep=2mm}

\newcommand{\level}[1]{\texttt{#1}}

\title{The Citation Provenance Chain and the Trust Registry:\\
\large How Clio Knows What She Knows, and How She Knows How Well She Knows It}
\author{Robin Langer \and Claude}
\date{4 July 2026}

\begin{document}
\maketitle

\begin{abstract}
Clio is an autonomous research agent living in a Docker container, working on
algebraic combinatorics across four kinds of session: \emph{browse} (reading
the literature through sub-agents), \emph{dream} (memory consolidation),
\emph{prove} (attacking one theorem), and \emph{lean} (formalising one
result). Each cycle transforms her knowledge, and each transformation can
corrupt it: sub-agents hallucinate paper contents, compression drops
citations, proofs absorb hearsay as fact. This document describes the two
systems deployed on 2026-07-04 to make such corruption mechanically visible:
a \emph{citation provenance chain} (a sources index with extraction levels
and a corrections history, plus a validator) and a \emph{proof-search trust
registry} (a tree of attempts with trust levels, a boundary rule, and a
source-ceiling rule). Both are advisory, stdlib-only, and grow lazily along
the paths actually walked. We explain what happens in each cycle, what can
go wrong, how old unverified claims are handled, and give worked examples,
including the incident that motivated the whole construction: a phantom
determinant that lived in her memory for two research cycles.
\end{abstract}

\tableofcontents

\section{The problem: knowledge laundering}

Clio never reads papers directly during browse sessions. Her main context is
scarce, so she orchestrates: four sub-agents search arXiv, MathOverflow, the
web, and Semantic Scholar citation graphs, and each returns short structured
summaries. Clio synthesises these into a daily reading log. This is
efficient, and it is also the root of the problem, because it creates a
pipeline in which errors are progressively stripped of the evidence that
they are errors:

\begin{enumerate}[nosep]
  \item \textbf{Browse.} A sub-agent paraphrases a paper. The paraphrase may
    be wrong. The reading log records the paraphrase with an honest citation,
    but nothing records that it \emph{is} a paraphrase.
  \item \textbf{Dream.} Consolidation compresses verbose notes into
    connections and topic files. Under compression pressure, arXiv
    identifiers and theorem numbers (``locators'') are the first casualties.
    The claim survives; its provenance does not.
  \item \textbf{Prove.} A proof session cites the now-unsourced claim as
    established literature. It has become load-bearing.
\end{enumerate}

We call this \emph{knowledge laundering}: each stage removes information
about where a claim came from and how reliable that channel was, so that by
the time the claim matters, it looks indistinguishable from a verified fact.

\subsection{The motivating incident: the phantom determinant}

On (narrative date) 2026-07-08, Clio's dream cycle produced a ``crown
jewel'' connection recording that the paper of Gatzweiler and Krattenthaler,
arXiv 2502.06032, contains an LGV (Lindstr\"om--Gessel--Viennot) determinant
relevant to her three-row program. On 2026-07-10 she happened to deep-read
the paper and found that it contains \emph{no determinant at all}. The claim
was a browse-agent extraction hallucination. It had been load-bearing in her
reasoning for two full cycles, and it was caught only by a voluntary
re-read. Nothing in the system could have caught it, because nothing
recorded that her knowledge of that paper had passed through a single lossy
paraphrase.

Two further errors surfaced during the audit that followed:

\begin{itemize}[nosep]
  \item \textbf{Di Francesco--Vu 2606.12796, over-claim.} Her notes read
    Theorem B.1's prefactor as giving an ``order law'' via a
    coincident-rapidity degeneration. The paper assumes pairwise-distinct
    variables and never takes the limit $x_i \to x_j$. The degeneration step
    is absent from the paper. It is hers to prove, not a literature import.
  \item \textbf{2605.13432, author misattribution.} Prior feeds attributed
    the paper to Gunna--Wheeler--Zinn-Justin; the actual authors are
    Gunna--Yeliussizov.
\end{itemize}

All three are now recorded as permanent \texttt{corrections} entries in her
sources index (Section~\ref{sec:sources}). The design principle throughout:
\emph{a wrong extraction is never deleted, only annotated}. Deleting it
would erase the evidence that this failure mode exists.

\section{The sources index}
\label{sec:sources}

The first component is a single file,
\path{memory/reading/sources.json}, format \texttt{clio-sources-v1}. Every
paper Clio cites gets one entry:

\begin{code}
"2606.12796": {
  "title": "Rank recursion for q-Whittaker and Macdonald operators",
  "authors": "Di Francesco, Vu",
  "year": 2026,
  "extraction": "deep-read",
  "read": ["reading/2026-07-06.md", "reading/2026-07-09.md",
           "reading/2026-07-10.md"],
  "locators": ["Thm B.1", "Rem 3.8", "Prop 2.2"],
  "corrections": [
    { "date": "2026-07-10",
      "note": "07-06/07-09 read Thm B.1's prefactor as giving the
       order law via coincident-rapidity degeneration. The paper
       assumes pairwise-distinct variables and never takes
       x_i -> x_j. The degeneration step is absent. It is mine
       to prove, not a literature import." } ]
}
\end{code}

\subsection{Extraction levels}

The key field is \texttt{extraction}: not \emph{what} the paper says but
\emph{how the knowledge of it was acquired}. Four levels, strictly ordered:

\begin{center}
\begin{tabular}{@{}llp{7.6cm}@{}}
\toprule
Level & Who read what & Typical failure mode \\
\midrule
\level{agent-summary} & a sub-agent read it; Clio has only the
  paraphrase & hallucinated contents (the phantom determinant),
  misattribution, over-compression \\
\level{abstract} & Clio read the abstract herself & the abstract
  oversells or omits the hypothesis that matters \\
\level{deep-read} & Clio read the full text & subtle over-claims
  (the missing degeneration step) \\
\level{verified-quote} & a specific statement copied verbatim and
  checked against the source & essentially none, for that statement \\
\bottomrule
\end{tabular}
\end{center}

A common misunderstanding is worth dispelling: \level{agent-summary} does
\emph{not} mean ``an old paper whose citation we lost.'' The citation is
fully known: arXiv ID, title, authors, and the exact reading logs it entered
through. What is untrusted is the \emph{content channel}. The paper is often
brand new; the knowledge of it is hearsay.

\subsection{Citation convention}

Inline, at the point of use, with a locator when a specific result is
invoked: \texttt{(arXiv:2606.12796, Thm B.1)}. Author names alone are not
citations for claims. This matters because the tooling extracts citations by
scanning for arXiv-ID-shaped tokens (a regex requiring a plausible month,
so decimals like $3.14159$ do not match); a claim ``by Di Francesco--Vu''
with no ID is invisible to every check described below.

\subsection{The validator}

\path{code/citation_check.py} (stdlib only; Clio's container has no PyYAML,
no pip installs) does four things:

\begin{itemize}[nosep]
  \item validates the index itself: known extraction levels, existing
    \texttt{read} files, well-formed corrections;
  \item checks that every arXiv ID cited in given markdown files resolves in
    the index, flagging \textbf{unregistered} sources;
  \item \texttt{--report footprint <file>}: for each source cited by a file,
    its extraction level and correction flags, plus the file's
    \emph{provenance floor}: the lowest extraction level among its
    citations. A connection with floor \level{agent-summary} is a connection
    that could be resting entirely on hearsay;
  \item \texttt{--report shallow}: every \level{agent-summary} source cited
    by at least one connections or topics file, ordered by how many files
    lean on it. This is the \emph{deep-read worklist}: the papers whose
    verification would raise the most floors.
\end{itemize}

Example footprint output on a note citing one deep-read source and one
unknown ID:

\begin{code}
connections/conn.md
  1111.22222  UNREGISTERED
  2606.12796  deep-read      Di Francesco, Vu  [1 correction(s)]
  provenance floor: deep-read
\end{code}

\section{The proof-search registry and the trust boundary}
\label{sec:registry}

The second component predates the citation chain by a few hours and indexes
a different thing: not papers, but \emph{her own attempts}.

\subsection{The fundamental unit: a node is an attempt}

One JSON file per conjecture; the file holds a tree. Each node is one
\emph{attempt}: sometimes a subgoal (``$c=3$: interior and boundary''),
sometimes a lemma (``Number Lemma: $16 \mid H$''), sometimes a strategy
(``LGV Gaussian-binomial determinant of the closed form''), sometimes the
corpse of a strategy (a dead end, with the reason it died). The node schema:

\begin{code}
{ "id": "unique-within-tree",
  "approach": "one line: what this attempt was",
  "trust": "speculative | computed | proved | lean-verified
            | dead-end | in-progress | unclassified",
  "file": "2026-06-18-c4-interior-Jstar-even.md",
  "lean": "optional_sorry_free_declaration",
  "reason": "required iff dead-end: why it died",
  "sources": ["2606.12796"],
  "children": [] }
\end{code}

A node's justification has exactly three kinds of evidence, each with its
own check:

\begin{center}
\begin{tabular}{@{}lll@{}}
\toprule
Field & Evidence type & Checked by \\
\midrule
\texttt{children} & internal, recursive & boundary rule (\S\ref{sec:boundary}) \\
\texttt{file}, \texttt{lean} & her own artifact & file existence; sorry-free build \\
\texttt{sources} & external literature & ceiling rule (\S\ref{sec:ceiling}) \\
\bottomrule
\end{tabular}
\end{center}

\subsection{Trust levels}

Ordered: \level{speculative} $<$ \level{computed} $<$ \level{proved} $<$
\level{lean-verified}. Outside the order: \level{dead-end} (abandoned; must
carry a \texttt{reason}), \level{in-progress} (open), and
\level{unclassified} (Section~\ref{sec:lazy}). The levels mean:

\begin{itemize}[nosep]
  \item \level{speculative}: an idea; no evidence beyond plausibility.
  \item \level{computed}: machine-checked evidence (verified for $n \le 10$,
    a counterexample found, a formula fitted), but no proof.
  \item \level{proved}: a written argument that holds, every step justified.
  \item \level{lean-verified}: a sorry-free Lean declaration, standard
    axioms only.
\end{itemize}

\subsection{The trust boundary}
\label{sec:boundary}

\begin{quote}
\textbf{Boundary rule.} A node may claim \level{proved} or
\level{lean-verified} only if every non-dead-end child is at least
\level{proved}.
\end{quote}

The picture: draw the tree and colour every node at \level{proved} or above.
The rule forces the coloured region to be \emph{upward closed along
justification}: a coloured node can only stand on coloured (or explicitly
abandoned) ground. The frontier between the coloured region and the rest is
the trust boundary. Above it, claims are justified by their entire subtree;
below it, everything is exploratory. The three report flags of the validator
are exactly views of this geometry: \texttt{successful-path} prints the
coloured skeleton, \texttt{frontier} prints the open nodes just below the
boundary, \texttt{dead-ends} prints the pruned branches with their reasons.

Worked example, from her real three-row registry. The $c=4$ subgoal:

\begin{code}
c4 [proved]
 |- c4-interior [proved]
 |    |- c4-number-lemma [proved]
 |    |- c4-linear-phi [dead-end]
 |         reason: v2(H) has a constant floor (>= 4),
 |                 not linear growth in phi
 |         |- c4-linear-phi-refutation [computed]
 |- c4-boundary [lean-verified]
\end{code}

\texttt{c4} may claim \level{proved}: its children are \level{proved} and
\level{lean-verified}, and the dead end does not count against it. Note the
refutation: a dead end's \texttt{reason} is itself a claim, and a wrong
reason silently prunes a live branch, which is the most expensive error in a
search. So dead ends can carry a \emph{child holding the refutation} at an
ordinary trust level. Here the refutation is \level{computed}: a
counterexample computation. A dead end with no refutation child was killed
on judgment; if she is ever tempted back down that route, the reason gets
re-checked first.

If she edited \texttt{c4-interior} down to \level{computed} while leaving
\texttt{c4} at \level{proved}, the validator would reject:

\begin{code}
root/c4: claims 'proved' but child 'c4-interior' is 'computed'
(boundary rule: non-dead-end children must be at least 'proved')
\end{code}

\subsection{External sources and the ceiling rule}
\label{sec:ceiling}

The boundary rule governs internal dependencies. The 2026-07-04 integration
extends the same idea to external ones:

\begin{quote}
\textbf{Ceiling rule.} A node claiming \level{proved} or
\level{lean-verified} whose cited sources are \emph{all} at
\level{agent-summary} is rejected: ``deep-read at least one before it is
load-bearing.''
\end{quote}

A browse-agent's paraphrase may be cited freely, exactly as an
\level{unclassified} child may be cited freely; what it cannot do is carry a
proof. Unknown IDs in a \texttt{sources} list produce warnings, not errors
(exit 0); a node with no \texttt{sources} field gets no check at all, which
keeps the rule backward compatible with registries written before the field
existed.

Live example. The node \texttt{lgv-det} in the three-row registry proposes
an LGV determinant route and notes it is ``the same determinant as the
Di Francesco--Vu order-law Cauchy det.'' Originally that dependency lived
only in prose, invisible to tooling. It now carries
\verb|"sources": ["2606.12796"]|. Today the node is \level{speculative}, so
the ceiling rule is silent. But the trap is armed: any future promotion to
\level{proved} happens in full view of that source's entry, including the
correction stating that the degeneration step \emph{is not in the paper}.
The node can no longer forget where its determinant came from.

Had the source been at \level{agent-summary}, promotion would fail outright:

\begin{code}
root/general-c/lgv-det: claims 'proved' but every cited source
is 'agent-summary' -- deep-read at least one before it is
load-bearing
\end{code}

\section{Old claims and lazy backfill}
\label{sec:lazy}

Clio has hundreds of pre-registry proof files and dozens of pre-index
reading logs. A full retroactive sweep would be both expensive and
epistemically fraudulent: labels assigned from memory are trust inflation.
The rule instead is \emph{lazy backfill}: nothing is retro-fitted; trust is
re-earned exactly when something new depends on it.

\subsection{\texttt{unclassified}: the ten-second IOU}

When a new attempt cites an old result not yet in the registry, the whole
obligation is one node:

\begin{code}
{ "id": "old-content-floor",
  "approach": "master valuation formula from June",
  "trust": "unclassified",
  "file": "2026-06-17-generalc-g0-content-floor.md",
  "children": [] }
\end{code}

\level{unclassified} means ``a claim exists on disk; nobody has re-checked
it.'' It ranks below \level{proved}, so the boundary rule does the rest: the
new node cannot claim \level{proved} while it leans on the IOU, and the
validator points at exactly which one. Files nothing depends on stay
unregistered forever, which is fine.

\subsection{Discharging: the citation lift}

To discharge an \level{unclassified} node, Clio re-reads the old proof and
promotes honestly: \level{computed}, \level{proved}, \level{dead-end} with a
reason (old results do turn out wrong; recording why is the point), or
\level{lean-verified}. Since 2026-07-04 the procedure has one more step,
because the old file's \emph{literature debts} must not vanish with the
label:

\begin{enumerate}[nosep]
  \item Run \texttt{citation\_check.py --report footprint} on the old file.
    It lists every arXiv ID in the text, each one's extraction level, and
    the provenance floor.
  \item Register anything unregistered in \path{sources.json} at its honest
    level. For a months-old browse cycle that is usually
    \level{agent-summary}: honesty, not pessimism.
  \item Lift the IDs into the node's \texttt{sources} field.
  \item Promote. The ceiling rule now applies: the node cannot reach
    \level{proved} if everything it imported is still hearsay, so the
    promotion forces deep-reads of precisely the papers the old argument
    leaned on.
\end{enumerate}

The invariant of the whole system is visible here: \emph{doubt can be moved
and sharpened, never silently dropped}. The coarse flag (``re-check this
whole file'') is converted into fine-grained ones (per-paper extraction
levels); the total amount of scepticism is conserved until evidence
discharges it.

One honest limitation. The lift is purely textual: the regex extracts IDs
that past-Clio wrote down. A claim she used silently, or cited by author
name only, leaves no machine-readable debt, and the node can reach
\level{proved} with no record that anything external was involved. The
provenance quality of old work is bounded by the citation discipline of the
Clio who wrote it. The compensating bet is temporal: the prompt hooks
(Section~\ref{sec:cycles}) make every file written from 2026-07-04 onward
cite IDs at the point of use, so extraction is lossless for new work and the
lossy cases are confined to a shrinking, quarantined corpus that is touched
only when something leans on it.

\section{The four cycles: what happens, what can go wrong}
\label{sec:cycles}

Each session type is a different epistemic mode, and each got hooks matched
to the claims born there. Discipline is installed twice over: as prose
instructions in the prompt, and, where possible, as a mechanical check the
session runs before ending.

\subsection{Browse: open-world acquisition}

\textbf{What happens.} Clio reads her seed file and recent memory, generates
compound keywords, and dispatches four agents in parallel (arXiv listings,
MathOverflow, general web, Semantic Scholar citation trails). When they
return she writes the daily reading log: papers, summaries, connections to
seed themes, open questions.

\textbf{What can go wrong.} This is where hearsay enters. Sub-agents
hallucinate contents (the phantom determinant), misattribute authors
(2605.13432), or summarise an abstract as if it were the paper. The log
itself cites correctly; the corruption is invisible because the log does not
say \emph{how} each fact was obtained.

\textbf{Hooks.} The synthesis step now includes: update
\path{sources.json} for every paper that got a real summary, at the honest
extraction level, noting which agent read it and whether it saw HTML full
text or only the abstract page; if today's reading contradicts an earlier
note, add a correction (never delete the wrong one). The session ends by
running \texttt{citation\_check.py} on the day's log: every ID cited must
resolve in the index, so a paper summarised but never registered is caught
the same day.

\subsection{Dream: compression under pressure}

\textbf{What happens.} A multi-pass streaming consolidation (her context is
small, so she reads batches, writes a checkpoint, and restarts). Five
phases: replay, associate, consolidate, prune, surface. The associate phase
births genuinely new claims: connections of the form ``X relates to Y
because Z, and this suggests Q,'' which she calls the crown jewels.

\textbf{What can go wrong.} Compression is exactly where locators die. A
note reading ``2502.06032 proves the crown identity via an LGV determinant
(Thm 4.2)'' becomes ``the G--K determinant gives the crown identity''
becomes ``the crown identity is known.'' Each rewrite is locally reasonable;
the composite is laundering. The phantom determinant survived two cycles
because consolidation carried the claim while shedding everything that would
have flagged it as hearsay.

\textbf{Hooks.} Consolidation preserves arXiv IDs and locators on every
external claim; compression never drops a citation; a connection resting on
an \level{agent-summary} source must say so inline, because that flag is
what tells future-Clio to deep-read before the claim becomes load-bearing.
Before exiting, she runs \texttt{citation\_check.py} on any connections or
topics files she rewrote, catching IDs orphaned by the rewrite. Residual
gap, stated honestly: a validator cannot flag the absence of an ID it never
saw, so a citation deleted from prose \emph{entirely} is still only guarded
by instruction. The \texttt{read} trail in \path{sources.json} at least
makes such a loss recoverable: the paper's entry records which logs it
entered through.

\subsection{Prove: closed-world reasoning}

\textbf{What happens.} One theorem, the whole session. No email, no
browsing: the network ban is deliberate, so that no fact can enter a proof
through a hasty, unrecorded read. Sub-agents may compute (check a lemma for
$n \le 10$, hunt counterexamples) but Clio holds the proof thread.

\textbf{What can go wrong.} Two things. Internally: a chain of lemmas where
one link is only \level{computed} evidence gets written up as a proof.
Externally: a claim from memory, laundered through a dream cycle, gets cited
as a known theorem. The missing-degeneration incident is the second failure
in its subtlest form: the paper is real, the theorem is real, and the single
step Clio actually needed is not in it.

\textbf{Hooks.} External facts cite ID and locator at the point of use and
go into the registry node's \texttt{sources} field. An \level{agent-summary}
source is not load-bearing: deep-read first (in a browse cycle; the shallow
report queues it) or mark the step as a gap. Pre-registry citations become
\level{unclassified} children. The session ends by running
\texttt{registry\_validate.py}: boundary rule, ceiling rule, dead-end
reasons, file existence. When she hits the ceiling, the effect is that ``I
need this paper'' becomes a queued, typed request rather than an impulsive
fetch: prove names the blocking source, the shallow worklist ranks it,
browse deep-reads it, and the next prove session finds the ceiling lifted.

\subsection{Lean: formalisation}

\textbf{What happens.} One target declaration, already proved on paper.
Small steps, frequent builds; \texttt{sorry} is a bookmark, not a result;
axioms restricted to the standard three.

\textbf{What can go wrong.} The paper proof has a gap that formalisation
exposes (the desired outcome, in fact), or the temptation arises to patch it
with new mathematics mid-session, or to mark a node \level{lean-verified}
while sorries remain.

\textbf{Hooks.} Citations from the paper proof are carried into the Lean
file as docstrings or comments on any declaration formalising an external
result, so even the formal artifact remembers its literature. A node is set
to \level{lean-verified} only when the declaration builds sorry-free, and
the validator run closes the session.

\subsection{Summary table}

\begin{center}
\small
\begin{tabular}{@{}lp{3.4cm}p{3.9cm}p{3.9cm}@{}}
\toprule
Cycle & New claims born & Main failure mode & Mechanical check at exit \\
\midrule
browse & ``paper X contains Y'' & agent hallucination,
  misattribution & \texttt{citation\_check} on the day's log \\
dream & connections, syntheses & citation loss under
  compression & \texttt{citation\_check} on rewritten files \\
prove & lemmas, registry nodes & hearsay becomes
  load-bearing & \texttt{registry\_validate}: boundary + ceiling \\
lean & formal declarations & premature
  \level{lean-verified} & \texttt{lake build}, sorry count, axiom print \\
\bottomrule
\end{tabular}
\end{center}

\section{Three kinds of ``we don't trust this yet''}

The system distinguishes three states that are easy to conflate:

\begin{description}[nosep]
  \item[Unregistered.] An arXiv ID appears in a note but has no entry in
    \path{sources.json}. \emph{Who is this?} Caught by
    \texttt{citation\_check.py}, exit 1.
  \item[\level{agent-summary}.] Provenance fully known, transmission lossy:
    someone else's work, seen through a summariser. \emph{We know who, but
    only by hearsay.} Caught by the ceiling rule at promotion time and
    surfaced by the shallow worklist.
  \item[\level{unclassified}.] Provenance known, verification stale: her own
    past claim, gathering dust. \emph{Our own work, unre-checked.} Caught by
    the boundary rule at promotion time; discharged by re-read plus citation
    lift.
\end{description}

The first is a bookkeeping error; the second and third are the external and
internal faces of the same epistemic state, and the two rules treat them
symmetrically: cite freely, but nothing above the boundary may rest on them.

\section{What the system does not do}

Honesty about scope, because the gaps define the next projects:

\begin{itemize}[nosep]
  \item \textbf{No online verification.} Nothing goes to the network to
    check a citation. Prove sessions are forbidden to browse by design;
    deep-reads happen in browse cycles, routed by the shallow worklist. The
    system verifies citations that exist in the text; it does not discover
    missing ones.
  \item \textbf{No completeness or priority search.} Nothing asks ``does
    this claim of hers already exist in the literature?'' or ``has someone
    published a contradiction?'' That would be a literature-search tool, not
    a validator.
  \item \textbf{Prose-only mentions are invisible.} ``By Di Francesco--Vu''
    without an ID escapes every check. The convention (IDs at point of use)
    exists precisely to shrink this hole for all new writing.
  \item \textbf{Advisory, not blocking.} Both validators print readable
    messages and exit nonzero on real violations, but no automation prevents
    Clio from ignoring them. This is deliberate: she is a researcher, not a
    pipeline stage, and the tools are instruments she reads, not gates she
    passes. The bet is that visible, precise, cheap checks get used; the
    first browse and prove cycles after deployment will show whether the bet
    pays.
\end{itemize}

\section{The mathematical reading}

The registry is a directed container in the sense of Ahman--Chapman--Uustalu:
shapes are search trees, positions are paths from the root. The structure
maps are exactly the operations a resumed session needs: \texttt{root} is
the conjecture (the empty path), \texttt{sub}$(t,p)$ is the sub-search at
path $p$ with dead ends included, and \texttt{shift}$(t,p,q)$ composes
paths, viewing $q$ inside the sub-search at $p$ globally. The induced
comonad has \texttt{extract} reading the root's trust (``is it proved?'')
and \texttt{duplicate} decorating every node with the full sub-search below
it, so that standing at any attempt, the complete history beneath is
available. The three report flags are coKleisli morphisms
$W(\mathrm{Registry}) \to \mathrm{Report}$: strategies that inspect the
whole tree of attempts, failures included, before summarising. This
structure is formalised, all five comonad laws and no \texttt{sorry}, in
\path{ProofSearch.lean} in the containers repository.

The citation chain has a reading in the same vocabulary: reading a paper is
pulling positions back along a container morphism from Clio's memory to the
literature, and the extraction level records how much of that morphism has
actually been traversed rather than presumed. It is the first instance of a
planned pattern of cross-registry morphisms, where trust in one structure is
constrained by typed references into another.

\section{Deployment facts}

For the record (all deployed 2026-07-04, real date; Clio's narrative clock
runs a few days ahead, so her file dates differ):

\begin{itemize}[nosep]
  \item \path{memory/reading/sources.json}: 19 sources backfilled from
    reading logs 07-06 through 07-10 at honest levels, 3 corrections. Around
    45 re-mentioned hub papers deliberately left out: lazy backfill applies
    to the index too.
  \item \path{memory/reading/CITATIONS-README.md}: conventions, addressed to
    Clio.
  \item \path{code/citation_check.py} and the extended
    \path{code/registry_validate.py}: both stdlib-only, both run clean in
    place against her real tree.
  \item Prompt hooks in \path{browse-prompt.md}, \path{dream-prompt.md},
    \path{DREAM.md}, \path{prove-prompt.md}, \path{lean-prompt.md}. Scripts
    are baked into her Docker image, so host edits were copied into the
    running container and verified byte-identical.
  \item The \texttt{lgv-det} node seeded with its Di Francesco--Vu source as
    the first live use of the \texttt{sources} field.
\end{itemize}

No Clio session has run since deployment. Adoption is the open question, and
the design's main mitigation for it: the hooks meet her at all four sites in
her own terse style, the tools are cheap and readable, and the conventions
formalise things she already half-did rather than imposing a foreign
workflow.

\end{document}
